\chapter*{Introduction Générale}
\addcontentsline{toc}{chapter}{Introduction Générale}
\markboth{Introduction Générale}{Introduction Générale}

\begin{spacing}{1.25}
{\large

Dans un contexte économique où la commande publique et les marchés d'entreprise
représentent des volumes financiers considérables, la gestion rigoureuse des appels
d'offres constitue un enjeu stratégique majeur. Pour les organisations qui publient
et évaluent régulièrement ces appels, la qualité du processus détermine directement
la compétitivité, la conformité réglementaire et la maîtrise des délais.

\todo[inline, color=orange!30]{À COMPLÉTER : contexte spécifique N+ONE — secteur datacenter, volume d'appels d'offres, positionnement marché}

\section*{Contexte du stage}

Ce stage de deuxième année a été effectué au sein de \textbf{N+ONE DATACENTERS},
société spécialisée dans l'hébergement et la gestion d'infrastructure datacenter,
basée à Casablanca, au Maroc. D'une durée d'un mois, du 6 août au 6 septembre 2025,
il s'inscrit dans le cadre de la formation en Génie Logiciel à l'École Nationale
Supérieure d'Informatique et d'Analyse des Systèmes (ENSIAS) de Rabat.

\todo[inline, color=orange!30]{À COMPLÉTER : description activité N+ONE, nombre d'employés, nombre de datacenters, clients types}

\section*{Problématique}

Avant ce projet, la gestion des appels d'offres chez N+ONE reposait sur des outils
non spécialisés : échanges par email, suivi sous tableurs Excel et stockage de documents
sur disques partagés. Cette approche fragmentée engendrait plusieurs problèmes
opérationnels :

\begin{itemize}
    \item \textbf{Absence de centralisation} : les informations étaient dispersées
    entre boîtes mail, fichiers Excel et répertoires réseau, rendant le suivi global
    impossible sans consolidation manuelle.

    \item \textbf{Risque de dépassement de délais} : aucun mécanisme d'alerte automatique
    ne prévenait les responsables à l'approche des dates limites de soumission
    ou d'évaluation.

    \item \textbf{Sécurité insuffisante} : les offres financières (données sensibles)
    circulaient par email ou étaient stockées en clair sur des partages réseau,
    sans chiffrement ni contrôle d'accès granulaire.

    \item \textbf{Traçabilité absente} : aucun historique ne permettait de savoir
    qui avait modifié quoi, ni quand, compromettant la conformité et la gouvernance.

    \item \textbf{Surcharge avant les échéances} : la réception de plusieurs soumissions
    simultanées en fin de période créait un goulot d'étranglement,
    les agents devant saisir manuellement les informations des candidats.
\end{itemize}

\section*{Objectifs du stage}

L'objectif principal était de concevoir et développer une \textbf{application web
sécurisée de gestion des appels d'offres}, couvrant l'intégralité du cycle de vie :
de la création de l'offre jusqu'à l'attribution du marché, en passant par la publication,
la collecte des soumissions et l'évaluation des candidats.

Les objectifs techniques et fonctionnels définis en début de stage étaient les suivants :

\begin{enumerate}
    \item Authentification sécurisée avec gestion des rôles (RBAC) ;
    \item Cycle de vie complet des appels d'offres avec transitions d'état contrôlées ;
    \item Portail public permettant aux sociétés externes de déposer leurs offres
    sans compte utilisateur ;
    \item Chiffrement AES-256-GCM des documents financiers sensibles ;
    \item Notifications automatiques (délais, changements de statut) ;
    \item Génération de rapports PDF ;
    \item Journal d'audit traçant toutes les actions sur les ressources ;
    \item Déploiement Dockerisé reproductible avec pipeline CI/CD.
\end{enumerate}

\section*{Plan du rapport}

Ce rapport s'organise en trois parties principales.

La \textbf{Partie 1} présente le cahier des charges : après une introduction à
l'organisme d'accueil et au domaine des appels d'offres, elle détaille l'analyse
des besoins fonctionnels et non-fonctionnels, l'étude des solutions existantes
et la justification du développement sur mesure.

La \textbf{Partie 2} couvre l'analyse, la conception et l'architecture : choix
technologiques argumentés, planification, modélisation UML (cas d'utilisation,
diagrammes d'activité, séquences) et conception détaillée du système (modèle de
données, API REST, architecture de sécurité).

La \textbf{Partie 3} décrit la réalisation et la mise en œuvre : implémentation
des modules backend et frontend, présentation des interfaces, et résultats
des tests de validation.

}
\end{spacing}

\newpage
